\chapter{Testing a Model on Your Own Task}
\label{ch:testing}

Published benchmarks measure someone else's task. A model that does well on
them can fail on a reworded question (\Cref{ch:landscape}), and the ecological
studies in \Cref{ch:applications} show how unevenly the same models perform
from one task to the next. Whether a model is good enough for your data is
therefore something to test on your data. This chapter collects ways in which
other projects have tested models, sets out a procedure for ecological tasks,
and ends with a small exercise in R that runs the procedure from start to
finish.

\section{More than one number}

A single accuracy score hides most of what matters. Four examples show other
ways to test.

\textbf{Test several things at once.} Karenina, from EMBL-EBI, lets experts
say in plain language what a sound answer needs and what behaviour to test,
and turns that into runnable checks. Its authors describe it in a preprint of
September 2026, not yet peer reviewed \parencite{carli2026karenina}. On 144
drug-discovery questions, live access to the Open Targets database raised the
average accuracy of seven models. Yet 11 per cent of the answers given with
that access were correct but unsupported by the retrieved evidence. Invented PubMed
identifiers gave way to real papers cited for claims they do not make, and
tool use cost many times more tokens. When a user confidently pushed back with
false information, models often abandoned correct answers: sycophancy,
measured (\Cref{ch:powers}). A second model acting as a guard flagged almost
all of those reversals. On an agent benchmark, switching the harness alone
changed one model's score, and reading the agents' saved code separated flawed
biology from failed computing.

\textbf{Ask what a benchmark can prove.} \textcite{saezrodriguez2026benchmarking}
argue that foundation models are hard to benchmark because they are meant to
embody the processes that produced the data. They ask whether such models can be
refuted, verified, or only judged by usefulness, and what part the scientific
community should play in testing them.

\textbf{Test against the previous version.} When EMBL-EBI rewrote InterProScan,
its protein annotation pipeline, as a Nextflow workflow, the authors compared
the new version's annotations of all Swiss-Prot proteins with the old one's. They
found near-identical precision and sensitivity, at about twice the speed on large
proteomes \parencite{blum2026interproscan}.

\textbf{Test the work as it is done.} OmicsBench, also a preprint, poses a
bioinformatics study to an agent as a chain of dependent steps, one turn at a
time, rather than as a single prompt \parencite{shaked2026omicsbench}. Agents
scored lower on steps that built on their own earlier answers. Its graders
match genes and organisms through reference databases, so a valid synonym is
not marked wrong. Of its 97 tasks, all but three are kept private, so that the
benchmark can still test models released later. Among the failures the
authors describe are agents that wrote answers from memory without opening the
data.

\section{A procedure for ecological tasks}

The ecological examples in \Cref{ch:applications} suggest a procedure of their
own. The following is my proposal for comparing extraction, classification and
study-retrieval workflows, drawing on their different successes and failures
\parencite{gougherty2024extraction,domazetoski2025traits,dorm2026ecologicalknowledge,munch2026phenotypes}
(\Cref{fig:procedure}).

\begin{enumerate}
  \item \textbf{Fix the task and the reference answers.} Specify the fields,
    trait definitions, units and acceptable evidence. Have two experts annotate
    a sample independently, report how often they agreed, resolve disagreements
    and retain uncertain cases as uncertain. Record the
    \term{reference-database version} that the key relies on, such as the
    taxonomic backbone or sequence database, for the key and for the system
    under test, and score again when either changes. Set aside a final test set
    before adjusting prompts, rules or thresholds, or writing a baseline.
    Keep descriptions of the same species, or samples from the same study,
    together when testing transfer to new species or studies. For medicine,
    \textcite{salaudeen2026failures} recommend testing on held-out hospitals,
    periods or patient groups, because models that score well where they were
    developed can fail elsewhere. For field notes, the equivalent is notes from
    another observer, site or season.
  \item \textbf{Compare like with like.} Run keyword matching or a simple
    classifier alongside the language model. Include the established domain
    method where one exists. Test unaided answers, answers given a fixed source
    passage, and answers using retrieval separately. Record which documents
    each system could access; these conditions test different abilities.
  \item \textbf{Score the errors that matter.} For extraction, report precision
    and recall for each field, plus how often missing information becomes an
    invented value. Check numerical errors after converting units, stating any
    transformation used. For traits, show each class's results and balanced
    accuracy, so common classes cannot conceal failure on rare ones. For names,
    report results by taxonomic rank as well as overall, since accuracy at
    genus level can hide failure at species level.
    Count malformed outputs, failures and unanswered cases across the full test
    set, even when calculating accuracy only on valid answers.
  \item \textbf{Test uncertainty and unfamiliar cases.} Include absent fields,
    rare taxa and difficult negative examples. Add controls for benchmark
    contamination (\Cref{ch:powers}), since public reference data may be in a
    model's training data. Invented names cannot have been memorised
    \parencite{munch2026phenotypes}, and models know least about the most
    specific ranks \parencite{sun2024taxoglimpse}. Report
    \term{coverage}, the proportion of cases receiving an answer, alongside
    accuracy. If confidence controls which answers are accepted, check whether
    higher confidence predicts fewer errors. Measure the errors and
    coverage again at each proposed threshold.
  \item \textbf{Evaluate retrieval and evidence separately.} For literature
    assistants, measure how many expert-selected relevant papers appear in the
    first few results. Then check whether each generated claim is supported by
    its cited passage. If a second model does that checking, compare its
    verdicts with an expert's on a sample first (\Cref{ch:localfirst}).
    For microbiome comparisons, add independent expert
    judgements of study compatibility. Agreement between text and biological
    profiles is useful corroboration, but does not substitute for those checks.
  \item \textbf{Measure variation and the whole cost.} Repeat runs and modest
    prompt changes, checking agreement on individual records. Report results
    by taxon, biome or source collection, with sample sizes and uncertainty
    intervals computed over the right unit, and compare systems on the same
    items (the exercise below shows how). Record runtime, token use, hardware
    and human correction time.
    Judge a replacement by verified output per unit of work and by unacceptable
    errors, with the acceptance criteria fixed before comparing models.
    Expect this evaluation to be the main cost. The curators who built agentic
    pipelines for knowledgebases report that writing the software was cheap
    beside rating its outputs, which still needs people
    \parencite{pombase2026curation}. They also found that experienced
    curators caught more flaws in generated ontology terms than newer ones.
    They propose keeping every rejected suggestion and every entity a person
    added, so that false positives and misses can be counted.
\end{enumerate}

\begin{figure}[htbp]
\centering
\fig{ch14-procedure}
\caption{The procedure in six steps, from fixing the task and its reference
answers to measuring variation and the whole cost.}
\label{fig:procedure}
\end{figure}

A benchmark's comparison method needs checking as well as the language model.
Where a score rests on a matching rule or a similarity measure, try the rule on
cases whose answer is known. The exercise below scores a free-text effort field
by accepting a phrase when either contains the other. In one test run, that
rule accepted ``Deep soil core'' for ``Deep soil core, 10 to 20 cm'' and so
lost the depth. Repeating a judgement consistently establishes repeatability;
an independent reference is needed to assess correctness.

\section{An exercise: field notes}

The \code{workshop/testing} folder applies the procedure to a task small
enough to finish in an afternoon. Invented field notes from the survey of
\Cref{ch:firstsession} each record a station, a date, a habitat and at most one
species with its count. The task is to turn each note into one row of a table
with nine fields. Ten notes, N01 to N10, are the \term{development set}, and
the baseline below was written with them in view. Fifteen more, N11 to N25, were
written after the baseline as the \term{held-out set}. The baseline runs on
them unchanged, and the rules were fixed before any model saw them. Both sets
include failed samples, rough and partial counts, places without a station
number, species searched for and not found, tentative identifications and
misspelt names. These are the absent fields and difficult negative examples of
step~4.

\textbf{Keep the records ecology needs.} An earlier version of the exercise had
five fields and did not record a species that was searched for and not found.
The rules now keep it, with \code{status} set to \code{absent}, the habitat
searched and the effort spent. A lost or untaken sample is \code{not\_sampled},
and a number that covers only the first tray is marked \code{partial} in
\code{count\_qualifier}. An abundance or occupancy analysis needs exactly this
information to tell an absence from a gap (\Cref{sec:fitformodel}). The fields
also expose the most serious error an extractor can make here, the
\term{false zero}. This is a count of 0 written for a sample that was lost or
never taken, and it enters an analysis as a true absence.

\textbf{Fix the rules first.} \code{RULES.md} defines each field before
anything is run. A count is recorded only when an exact number is stated, in
digits or words, and it is never 0: ``roughly a dozen'' leaves it empty, with
\code{count\_qualifier} set to \code{approximate}. Names are copied as written,
misspellings and ``cf.'' included, because checking them is a separate step.
Habitats come from a closed list of six. Before looking at the keys in
\code{key/}, annotate the development notes yourself, ideally with a colleague
working separately, and note how often the two of you agree. Where you
disagree with the key, decide whether the rule or your reading was at fault,
and change the rules only now, before any model is tested.

\textbf{Run a baseline.} \code{R/01\_baseline.R} uses patterns and keywords
only. It has a regular expression for the station code, three date formats, an
exact match against the checklist and the number nearest the name, with keyword
lists for status, effort and the two qualifiers. It needs no model and gives
the same answer every time. Its patterns were written while looking at the
development notes, so its score there measures how well it fits them. Only the held-out
notes show how it does on notes it was not written for. The script also sets
the time locale to C, because month names such as ``May'' are otherwise read in
the computer's language. Without that line, it loses every long-form date on a
Greek or German system.

\textbf{Run a model.} \code{R/02\_extract.R} sends each note, in a fresh
conversation, to a local model through Ollama, using the \code{ellmer} package.
It asks for output that fits a fixed structure:

\begin{lstlisting}[language=R,caption={Sampling settings and the structure the model must fill, abridged from \code{R/02\_extract.R}}]
sampling <- list(temperature = 0.6, top_p = 0.95, top_k = 20)
p <- params(temperature = sampling$temperature, top_p = sampling$top_p,
            top_k = sampling$top_k, seed = if (is.na(seed)) NULL else seed)

type_note <- type_object(
  station = type_string("Station code S1 to S5, as written",
                        required = FALSE),
  date    = type_string("Sampling date as YYYY-MM-DD", required = FALSE),
  species = type_string("Scientific name exactly as written",
                        required = FALSE),
  count   = type_integer("Exact number of individuals", required = FALSE),
  habitat = type_enum(c("litter", "soil", "moss", "bark", "compost", "cave"),
                      required = FALSE),
  # ... status, effort, count_qualifier and id_qualifier: see the script
)
chat <- chat_ollama(
  model = model,
  system_prompt = paste("Extract one record from a field note.", rules, sep = "\n\n"),
  params = p
)
c1  <- chat$clone()                      # a fresh conversation for each note
rec <- c1$chat_structured(notes$text[i], type = type_note)
\end{lstlisting}

Each field is optional, so the model is allowed to leave it empty. Without
that, a model must fill every field and is pushed towards inventing values.
\code{status} and \code{count\_qualifier} are closed lists, like
\code{habitat}; \code{effort} and \code{id\_qualifier} are free text. Every run
uses the same sampling settings, taken from the model's documentation.
\textcite{miller2024errorbars} advises against lowering the temperature to
steady the scores, because that changes the model under test. An optional seed
makes the sampling repeatable. The script keeps every raw response, including
the reasoning that Qwen3 writes before its answer, and records the settings and
timings beside the output. Its arguments are the run's label, the set and, if
wanted, the seed. From inside \code{workshop/testing}:

\begin{lstlisting}[style=shell,caption={Baseline, three free runs, two seeded runs and a score, for each set}]
for set in dev heldout; do
  Rscript R/01_baseline.R $set
  for run in run1 run2 run3; do Rscript R/02_extract.R $run $set; done
  for run in fixedA fixedB; do Rscript R/02_extract.R $run $set 1; done
  Rscript R/03_score.R $set
done
\end{lstlisting}

\textbf{Score the errors that matter.} \code{R/03\_score.R} compares every file
in \code{results/<set>/} with that set's key, field by field. It separates
wrong values, missed values and invented values, where the key is empty but the
output is not (\Cref{fig:scoring}). Every field must match exactly except
\code{effort}, which counts as correct when either phrase contains the other.
The script then prints each disagreement, so that it can be read, not only
counted.

\begin{figure}[htbp]
\centering
\fig{ch14-scoring-grid}
\caption{Scoring against the key. Wrong, missed and invented values are
different failures. An invented value, where the key is empty and the output
is not, corrupts a database most quietly; a count of 0 for a lost sample, a
false zero, is one.}
\label{fig:scoring}
\end{figure}

When the exercise was tested for this book, on 2 November 2026, the baseline
got 83 of the 90 development cells right and 98 of the 135 held-out cells, the
same under a Greek locale. Qwen3 8B, a local model with 8 billion parameters,
scored 72 to 76 on the development notes, below the baseline in every run. On
the held-out notes it scored 107 to 115, above the baseline in every run.
There, the baseline missed dates in formats it had not seen, counts written as
words and names that did not match its checklist exactly. Twice it took the
wrong number as the count. The model's errors were of another kind. In every
run it wrote a count of 0 for a lost or unsampled note, a false zero, and a
habitat for a sample that was never taken. It also left a station code empty in
up to three notes per run. Its raw responses show that each time, its
reasoning had named the right station; the field was lost when the answer was
written out. The two seeded runs gave identical output, so the runs became
repeatable without becoming more correct. A run took about five minutes for ten
notes and eleven to thirteen for fifteen, on a laptop with other work running.

\textbf{Count notes, not cells.} The nine cells of a note are not independent:
a note that is hard to read tends to produce several wrong cells at once. The
\term{unit of replication} is therefore the note, and with ten or fifteen notes
any interval is wide. \code{03\_score.R} prints two 95 per cent intervals for
each run. The exact binomial (Clopper--Pearson) interval from
\code{binom.test()} treats every cell as independent. The bootstrap interval resamples whole notes, so it
respects the unit, and it is often wider. For the seeded runs' 115 of 135
held-out cells, they run from 0.78 to 0.91 and from 0.75 to 0.95. Ignoring
such grouping makes an evaluation look more precise than it is
\parencite{miller2024errorbars}. With tens of items, \textcite{bowyer2025clt}
find that intervals from the bootstrap and from the textbook normal
approximation are often too narrow, while the exact binomial interval errs on
the wide side. They recommend the Wilson interval, which
\code{prop.test(x, n, correct = FALSE)} gives, or a Bayesian one. For items in
groups, only their Bayesian method that models the groups kept its stated
coverage.

\textbf{Compare on the same notes.} When two systems are scored on the same
items, compare them item by item, because some items are easy or hard for both
\parencite{miller2024errorbars}. McNemar's test, \code{mcnemar.test()} in R,
does this for right-or-wrong scores. It sets aside the cells that both systems
got right or both got wrong, and asks whether the remaining disagreements lean
one way. On the held-out notes, the seeded runs were right where the baseline
was wrong in 34 cells and the reverse in 17, with $p = 0.025$. The three free
runs gave 0.27, 0.20 and 0.034. The script runs the test on cells, so it
shares the binomial interval's optimism, and I would not read $p = 0.025$ from
fifteen notes as a finding.

\textbf{Decide how many items you need.} The binomial arithmetic is simple. To
measure an error rate near 5 per cent to within 2 percentage points takes about
460 items ($1.96^2 \times 0.05 \times 0.95 / 0.02^2$), and even then
\code{binom.test(23, 460)} gives 3.2 to 7.4 per cent. Telling two systems apart
takes more. In a worked example with assumptions he calls reasonable,
\textcite{miller2024errorbars} needs about a thousand independent items to
detect a difference of 3 percentage points with 80 per cent power. Repeated runs belong in the same
account: the three free runs on the held-out notes ranged from 107 to 114.
Miller suggests answering each item several times and averaging its scores,
and OmicsBench averages each task's runs before computing any statistic
\parencite{shaked2026omicsbench}. Ten or fifteen notes can show what kinds of
error a system makes, but not which system makes fewer.

These numbers describe one small model, twenty-five notes and one afternoon,
and your runs will differ. They do not show that the model is better than the
baseline: the intervals overlap, and the paired test counted cells. What they
show is narrower. A baseline tuned on the notes it was scored on beat the model
there and fell behind it on notes written later, so a score on development
items says little about the next ones. The disagreements say more than the
totals. The baseline failed on formats it had not seen; the model failed at
leaving a count empty, the error that would corrupt an abundance table most
quietly. Keeping the patterns where they never failed, such as station codes,
and using the model elsewhere might beat either. That is a hypothesis for the
next held-out set, not a result of this one.

\textbf{Take it further.} Change \code{03\_score.R} so that the paired test, or
a bootstrap of the differences, works on notes rather than cells, and see how
much of the held-out difference survives. The model sometimes wrote the word
``empty'' into \code{id\_qualifier}, which is free text; make it a closed list,
as \code{status} is. Try a larger model, or a hosted one through another of
\code{ellmer}'s chat functions, and add it to the same table. Record the time
each run takes and how long you spend checking the output, as step~6 asks. The
\code{vitals} package does at scale what \code{03\_score.R} does
(\Cref{app:tutorials}). Then write twenty notes of your own, from your own
field books, and keep them aside as the test set for the next model you are
offered. Once anything has been tuned on them, they are development notes, and
the next test needs new ones.

The same discipline applies to every choice in \Cref{ch:localfirst,ch:hardware}.
A quantised model, a fine-tuned one or a new release earns its place by passing
the tests the previous one passed, judged against criteria fixed in advance. It
keeps that place only while a sample of its output keeps passing. In a review
of AI in health care, \textcite{salaudeen2026failures} describe an AI scribe
whose accuracy in assigning diagnosis codes fell from 79 to 35 per cent after a
change to note templates. Systematic monitoring caught the drop. When a model's
output becomes the data, \Cref{sec:outputasdata} describes the validation
sample that carries this check into the analysis.
